% Source: https://artofproblemsolving.com/community/c3h2519196_analytical_geometry
% Stable identity: aops-topic-2519196-q01
\textbf{(a)} We can write
\[P_k=\left(a+\dfrac kn(c-a),\ b+\dfrac kn(d-b)\right).\]
So $P_k$ is reached from $A$ by moving $\dfrac kn$ of the horizontal change $c-a$ and $\dfrac kn$ of the vertical change $d-b$. Hence consecutive points differ by the same step $\left(\dfrac{c-a}{n},\dfrac{d-b}{n}\right)$. By the distance formula,
\[P_kP_{k+1}=\sqrt{\left(\dfrac{c-a}{n}\right)^2+\left(\dfrac{d-b}{n}\right)^2}=\dfrac1n\sqrt{(c-a)^2+(d-b)^2}=\dfrac{AB}{n}.\]
All the steps are the same vector, so $P_0=A,P_1,\dots,P_n=B$ lie on one line, in this order, and the $n$ pieces are equal. (This is the midpoint formula for $n=2$.) For $n=3$, $k=1,2$ gives the stated trisection points.

\begin{asy}
usepackage("amsmath");
size(5.5cm);
defaultpen(fontsize(10pt));
pair A=(0,0), B=(6,4.5);
pair P=A+(B-A)/3, Q=A+2*(B-A)/3;
draw(A--B,linewidth(0.9));
draw(A--(P.x,A.y)--P,dashed+blue); draw(P--(Q.x,P.y)--Q,dashed+blue); draw(Q--(B.x,Q.y)--B,dashed+blue);
dot(A);dot(B);dot(P);dot(Q);
label("$A(a,b)$",A,W); label("$B(c,d)$",B,E);
label("$P_1$",P,NW); label("$P_2$",Q,NW);
label("$\dfrac{c-a}{3}$",(A.x+P.x)/2,S,blue+fontsize(8pt));
label("$\dfrac{d-b}{3}$",(P.x,(A.y+P.y)/2),E,blue+fontsize(8pt));
\end{asy}

The sketch shows $n=3$: the three small right triangles with legs $\dfrac{c-a}{3}$ and $\dfrac{d-b}{3}$ are congruent, which is the same argument in pictures.

\textbf{(b)} With $(a,b)=(-2,5)$ and $(c,d)=(7,-1)$:
\[\left(\dfrac{2(-2)+7}{3},\dfrac{2\cdot5+(-1)}{3}\right)=(1,3),\qquad\left(\dfrac{-2+2\cdot7}{3},\dfrac{5+2(-1)}{3}\right)=(4,1).\]
The trisection points are $\boxed{(1,3)\text{ and }(4,1)}$.
